\section*{Chapter \ref{ch:b3:surfaces-catalysis} --- Surfaces, Adsorption and Heterogeneous Catalysis}

\begin{solution}{exo:b3:surfaces-catalysis:1}
$K = \theta/((1 - \theta)p) = 0.40/(0.60 \times 2.0) = \qty{0.33}{kPa^{-1}}$; at \qty{10}{kPa}, $\theta = 3.33/4.33 = 0.77$.
\end{solution}

\begin{solution}{exo:b3:surfaces-catalysis:2}
$-20$: \omterm{def:b3:surfaces-catalysis:physi-chemi}{physisorption} (the order of a condensation enthalpy); $-150$: \omterm{def:b3:surfaces-catalysis:physi-chemi}{chemisorption}, which
alone can be dissociative.
\end{solution}

\begin{solution}{exo:b3:surfaces-catalysis:3}
(100): one atop, two bridge (four bonds shared by two atoms each), one four-fold
hollow (four hollows shared by four atoms). (111): one atop, three bridge, two
three-fold hollows (six triangles around each atom, each shared by three).
\end{solution}

\begin{solution}{exo:b3:surfaces-catalysis:4}
$\num{3.0e-6}/\num{1.5e-5} = \qty{0.20}{s^{-1}}$.
\end{solution}

\begin{solution}{exo:b3:surfaces-catalysis:5}
$\theta/(1 - \theta) = (Kp)^{1/2}$: $0.111$ at \qty{1.0}{kPa}, doubled at \qty{4.0}{kPa}: $0.222$, so $\theta = 0.18$.
\end{solution}

\begin{solution}{exo:b3:surfaces-catalysis:6}
Denominator $1 + 10 + 2.5 = 13.5$: $\theta_{\ce{CO}} = 0.74$, $\theta_{\ce{O2}} = 0.19$, free sites 0.07. The
Langmuir--Hinshelwood rate, proportional to $\theta_{\ce{CO}}\theta_{\ce{O2}}$, is limited by the
scarcity of oxygen on a surface crowded with CO.
\end{solution}

\begin{solution}{exo:b3:surfaces-catalysis:7}
$\Delta_{\mathrm{ads}}H = -R\ln(8.0/1.0)/(1/300 - 1/350) = -8.314 \times 2.079/\num{4.76e-4} = \qty{-36}{kJ/mol}$.
\end{solution}

\begin{solution}{exo:b3:surfaces-catalysis:8}
$n_{\mathrm m} = 120/22\,414 = \qty{5.35e-3}{mol}$; area $= \num{5.35e-3} \times \num{6.022e23} \times \num{0.162e-18} =
\qty{522}{m^2.g^{-1}}$.
\end{solution}

\begin{solution}{exo:b3:surfaces-catalysis:9}
Langmuir--Hinshelwood: both reactants must be adsorbed on neighbouring sites;
CO binds strongly and, at high pressure, leaves no room for oxygen.
\end{solution}

\begin{solution}{exo:b3:surfaces-catalysis:10}
Low coverage: $100 - 60 = \qty{40}{kJ/mol}$; high coverage: \qty{100}{kJ/mol}. The plot of $\ln r$
against $1/T$ is steep (slope $-100/R$) at low temperature, where the surface is
full, and flatter (slope $-40/R$) at high temperature, where it empties: a curve
bending between the two.
\end{solution}

\begin{solution}{exo:b3:surfaces-catalysis:11}
$4 \times 0.10 = 0.40$ of the sites are blocked: 60\,\% of the activity remains for one
site, about $0.60^2 = 36\,\%$ for two adjacent free sites.
\end{solution}

\begin{solution}{exo:b3:surfaces-catalysis:12}
The middle one (\qty{-250}{kJ/mol}): the first binds nitrogen too weakly to break \ce{N2}
readily, the third too strongly, so that nitrogen atoms stay on the surface
and block it.
\end{solution}

\begin{solution}{pb:b3:surfaces-catalysis:1}
\textbf{1.} Nitrogen condenses near \qty{77}{K} at atmospheric pressure: a full range of
relative pressures is reached, and multilayers form.
\textbf{2.} Below 0.05 the surface heterogeneity, above 0.30 condensation in the
pores, violate the BET assumptions.
\textbf{3.} \num{1.278e-3}, \num{2.410e-3}, \num{3.453e-3}, \num{4.621e-3}, \num{5.659e-3}, \num{6.813e-3} \unit{g.cm^{-3}}.
\textbf{4.} $v_{\mathrm m} = 1/(0.02205 + 0.000180) = \qty{45.0}{cm^3.g^{-1}}$.
\textbf{5.} $c = 1 + 0.02205/0.000180 = 124$: the first layer binds much more strongly than
the liquid does.
\textbf{6.} It is tiny compared with the slope times $x$, so its relative error is large;
but $v_{\mathrm m}$ depends on the sum of slope and intercept, dominated by the slope.
\textbf{7.} $45.0/22\,414 = \qty{2.01e-3}{mol.g^{-1}}$.
\textbf{8.} $\num{2.01e-3} \times \num{6.022e23} \times \num{0.162e-18} = \qty{196}{m^2.g^{-1}}$.
\textbf{9.} The alumina support: the exposed platinum (part III) has about \qty{0.9}{m^2.g^{-1}}.
\textbf{10.} Pores of molecular width fill completely at very low pressure, not layer
by layer, so no \omterm{def:b3:surfaces-catalysis:coverage}{monolayer} can be identified.
\textbf{11.} $2 \times 0.205/22\,414 = \qty{1.83e-5}{mol.g^{-1}}$.
\textbf{12.} $0.0100/195.08 = \qty{5.13e-5}{mol.g^{-1}}$.
\textbf{13.} $D = 1.83/5.13 = 0.357$.
\textbf{14.} $a^3/4 = (0.3923)^3/4 = \qty{0.01509}{nm^3}$.
\textbf{15.} $\frac{\sqrt3}{4}(0.3923)^2 = \qty{0.0666}{nm^2}$, the densest face; (100) and (110) give 0.0770 and
\qty{0.1088}{nm^2}, and the mean of the three site densities gives \qty{0.0807}{nm^2}.
\textbf{16.} A sphere holds $\pi d^3/6v_{\mathrm{at}}$ atoms, of which $\pi d^2/a_{\mathrm{at}}$ are on the surface:
$D = 6v_{\mathrm{at}}/(a_{\mathrm{at}}d)$.
\textbf{17.} $d = 6 \times 0.01509/(0.0807 \times 0.357) = \qty{3.1}{nm}$.
\textbf{18.} One H per surface Pt; spherical particles of one size (the result is a
surface-weighted mean); all the platinum reduced and accessible; no hydrogen
spilling over onto the support.
\textbf{19.} $\num{2.0e-5}/\num{1.83e-5} = \qty{1.1}{s^{-1}}$.
\textbf{20.} $\num{2.0e-5}/0.0100 = \qty{2.0e-3}{mol.g^{-1}.s^{-1}}$ per gram of platinum.
\textbf{21.} $D = 1.12/10 = 0.112$ instead of 0.357: a factor 3.2.
\textbf{22.} That the reaction is structure-sensitive: its \omterm{def:b3:complex-kinetics:enzyme}{active sites} (edges,
corners, particular faces) are not a constant fraction of the surface.
\textbf{23.} It counts the events per \omterm{def:b3:complex-kinetics:enzyme}{active site}, removing the effect of how much
metal is exposed.
\textbf{24.} \textbf{Mean platinum particle diameter $\approx \qty{3.1}{nm}$.}
\end{solution}
